\changecolours{ScoutGreen}
\section{Conclusão \& Próximos Passos}

% ------------------------------------------------------------------------------
\twocolframe[title={Síntese da Aula \& Orientações},leftcol=7cm,rightcol=7cm]{%
  \alert{Conceitos Consolidados}\par
  \itemseps[topsepi=2mm,itemsepi=3mm,labelsep=2mm,leftmargini=4mm]
  \begin{itemize}
    \item \textbf{Três Níveis ANSI/SPARC:} Nível Externo (visões particulares), Conceitual (modelo global de negócio) e Interno (armazenamento e índices em disco).
    \item \textbf{Independência de Dados:} Imunidade garantida por mapeamentos entre camadas, dividida em Lógica (visões e esquemas) e Física (arquivos e índices).
    \item \textbf{Dicionário de Metadados:} Auto-descrição do SGBD que viabiliza compilação segura, controle de permissões e otimização de custo (CBO).
    \item \textbf{Linguagem SQL Completa:} Domínio taxonômico das famílias DDL, DML, DQL, DCL e TCL.
  \end{itemize}
}{%
  \alert{Próximas Atividades (Cronograma)}\par
  \itemseps[topsepi=2mm,itemsepi=3mm,labelsep=2mm,leftmargini=4mm]
  \begin{itemize}
    \item \textbf{Aula 03 (06/10 - Terça-feira):} Fases do Projeto de Banco de Dados e Introdução ao Modelo Entidade-Relacionamento (MER).
    \item \textbf{Aula 04 (10/10 - Sábado Letivo):} Resolução e entrega assíncrona da \textbf{Lista de Exercícios 01 (N1)}, cobrindo os conceitos de SGBD, arquitetura e modelagem básica.
    \item \textbf{Canal de Mentoria:} Dúvidas sobre os scripts e exercícios podem ser enviadas diretamente via repositório ou atendimento presencial no laboratório.
  \end{itemize}
}
